% ================================================================
% Cover Letter — gear_fatigue_v2.7
% Target: International Journal of Fatigue (Elsevier)
% ================================================================
\documentclass[12pt,a4paper]{article}
\usepackage[margin=2.5cm]{geometry}
\usepackage{hyperref}

\begin{document}

\thispagestyle{empty}

\vspace*{1cm}
\hfill\begin{minipage}{0.45\textwidth}
\raggedleft
Zhilei Chen\\
Guangdong Peizheng College,\\
Data and Computer Science,\\
53 Peizheng Rd., Chini Town,\\
Huadu, Guangzhou,\\
Guangdong 510830, China\\[4pt]
\texttt{2604513@peizheng.edu.cn}
\end{minipage}

\vspace{1cm}
\noindent The Editor\\
International Journal of Fatigue\\
Elsevier

\vspace{0.8cm}
\noindent\textbf{Date:} 23 July 2026

\vspace{0.3cm}
\noindent\textbf{Re:} Submission of ``Quantifying Method-Induced Uncertainty in
Mechanical Reliability: A Factorial Variance-Decomposition Framework
Demonstrated on Gear Bending Fatigue''

\vspace{0.8cm}
\noindent Dear Editor,

\vspace{0.3cm}
I am pleased to submit the above manuscript for consideration for
publication in the \textit{International Journal of Fatigue}.

\vspace{0.3cm}
\noindent\textbf{What this paper does.}
When three engineers compute the bending fatigue life of the same spur
gear---same geometry, same material, same load---their predictions can differ
by four orders of magnitude. The divergence comes from methodological choices:
which standard to use (ISO~6336, AGMA~2001, or FKM), how to correct for mean
stress, what S--N curve to trust, how to penalize surface roughness. Each
choice, studied in isolation by prior work, has been shown to matter.
What has never been done---and what we present here---is to place all
of these choices into a single controlled factorial experiment and decompose
the total log-variance into its constituent sources, ranked by magnitude.

Using a full-factorial ANOVA over 144 combinations (5 switches, 116
finite-life entries), we find that two factors account for 62.0\% of total
variance: the choice of standard (31.1\%) and the choice of mean-stress
correction (30.9\%). Surface roughness (8.5\%) and a Standard~$\times$~Rz
interaction (5.0\%) are secondary. The S--N curve source and cumulative
damage rule---factors that the literature often treats as consequential---are
negligible ($<$1\% each). Bootstrap validation with 2000 draws confirms
stable rankings. A multi-torque comparison reveals that the factor ranking
is operating-point-dependent, not a universal constant.

All formulas have been verified against the original standards (ISO~6336-3:2019,
AGMA~2001-D04, FKM~7th ed.). S--N parameters are from the publicly
accessible Waterloo SMDIdbase. The analysis pipeline is released as an
open-source, auditable toolkit.

\vspace{0.3cm}
\noindent\textbf{How this differs from prior work---including recent
contributions in this journal.}
We are aware that the \textit{International Journal of Fatigue} has recently
published important work on standard-to-standard divergence in gear fatigue
assessment, notably Bonaiti et al.\ (\textit{Int.\ J.\ Fatigue}, Vol.~181, 2024,
108145), who compared ISO~6336 and AGMA~2001-D04 predictions against pulsator
test data. That study provides an invaluable experimental benchmark. Our
manuscript complements and extends it in three ways:

\begin{enumerate}
  \item \textbf{From binary comparison to factorial decomposition.}
        Bonaiti et al.\ ask whether ISO and AGMA differ (yes, they do). We
        ask: across all major methodological switches an engineer faces, how
        is the total log-variance partitioned---and which switch contributes
        the most? The factorial framework reveals that mean-stress correction
        (not examined by Bonaiti et al.) is statistically indistinguishable
        from standard choice as the largest variance source.

  \item \textbf{From two dimensions to five.} Bonaiti et al.\ examine
        standard choice and statistical elaboration method. We simultaneously
        vary five independent choices (standard, mean-stress correction,
        roughness grade, S--N source, damage rule), capturing interactions
        that binary comparisons miss.

  \item \textbf{From qualitative conclusion to quantitative ranking.}
        ``Standards differ'' is a qualitative statement. ``The standard
        accounts for 31.1\% of variance with a bootstrap S/N ratio of 5,
        while mean-stress correction accounts for 30.9\% with S/N of 4''
        is a quantitative one that enables engineering resource allocation:
        which uncertainty sources warrant the most modeling effort?

  \item \textbf{Bootstrap noise floor.} No prior gear fatigue UQ study has
        employed bootstrap resampling to distinguish genuine signal from
        finite-sample fluctuation, making our uncertainty budget a
        repeatable measurement rather than a one-shot estimate.
\end{enumerate}

The two studies are complementary, not competing: Bonaiti et al.\ provide
independent experimental evidence that the standard-choice effect exists; we
measure its magnitude relative to four other uncertainty sources and provide
a reusable framework applicable to any mechanical reliability calculation
where multiple handbooks or methods compete.

\vspace{0.3cm}
\noindent\textbf{Fit with the journal.}
The \textit{International Journal of Fatigue} is the natural home for this
work. The journal has published gear bending fatigue studies (Glodez et al.,
2002; Bonaiti et al., 2024), cumulative damage surveys (Fatemi \& Yang, 1998),
and mean-stress methodology papers---each examining one factor in
isolation. Our manuscript synthesizes these threads within a unified UQ
framework, and we explicitly cite and build upon the journal's own
published record. The methodology is not limited to gears: it applies to
any mechanical component where multiple standards or correction methods
compete.

\vspace{0.3cm}
\noindent\textbf{Additional information.}
This manuscript has not been published previously and is not under
consideration elsewhere. All authors have approved the submission. The
authors declare no competing financial interests.

\vspace{1cm}
\noindent Sincerely,

\vspace{1.5cm}
\noindent Zhilei Chen

\end{document}
